% ==============================================================================
% ПАРТИЯ B03-013: DISS-005-B0550 .. DISS-005-B0586
% Глава 1: Отпадание адаптивного РНМ (вопрос подбора К отс)
% Замечания: DISS-005-C0057 .. DISS-005-C0062
% Иллюстрации: image23.png, image24.png, image25.png, image26.png
% ==============================================================================

% DISS-005-B0550 .. DISS-005-B0553 (Примечание к таблице 1)
\textbf{Примечание:}
\begin{itemize}[leftmargin=*,itemsep=0pt]
    \item[+] ПЧ остался в работе после срабатывания БАВР
    \item[-] ПЧ отключился после срабатывания БАВР
    \item[\%] \% уровня колебаний считается от амплитудного значения колебания
\end{itemize}

% DISS-005-B0555
\subsection{Отпадание адаптивного РНМ (вопрос подбора К отс)}
\label{sec:ch1_adaptive_dropout}

% DISS-005-B0556
В момент завершения пуска мощных синхронных двигателей может наблюдаться колебательное затухание амплитуды тока, что видно из рисунка 1. Причём это колебание начинает разрастаться как по амплитуде, так и по времени (рис. 2), то есть отстроиться по времени не \commentnote{DISS-005-C0057}{Евдаков Алексей}{Было вот так, но ввиду необходимости неразглашения тайн я сократил описнаие.\newline\newline«В момент времени 1~с отношение тока, например 200~мс назад (2040~А) и нынешнего тока (1630~А) опустилось ниже границы в 0,8, и произошло «отпускание» РНМ.»}{получится}. Ввиду данных колебаний ток опустился ниже автоматически рассчитываемого программой порога определяемого уставками 200~мс и 0,8~о.е. (и рядом иных скрытых параметров) и произошло «отпускание» РНМ. При этом, ввиду пуска двигателя, напряжение опустилось приблизительно до 86\% и уже произошёл пуск по 27/27S (ЗМН). Следовательно, сразу же после пропажи блокировки от РНМ отработал БАВР.

% DISS-005-B0557
Для решения этой проблемы коэффициент отстройки был уменьшен с 0,8 до 0,6, и дана рекомендация, при пуске двигателя выводить ключ БАВР.

% DISS-005-B0559 .. DISS-005-B0562 (DISS-005-FIG-020 / DISS-005-TAB-014)
\begin{figure}[htbp]
    \centering
    \includegraphics[width=0.85\textwidth]{image23.png}
    \caption[Рисунок 1. Общий вид осциллограммы]{Рисунок 1. Общий вид осциллограммы.\newline{\small Данные файлы находятся в папке «\_вспомогательные файлы\textbackslash{}ДКС-1\textbackslash{}Отпадание адаптивного РНМ при пуске мощного \commentnote{DISS-005-C0058}{Евдаков Алексей}{t00050\_210906\_0606\_\_0008}{\commentnote{DISS-005-C0059}{Евдаков Алексей}{На «ДКС-1 ЗРУ-10кВ СШ2-4» в 06:06 (06.09.2021) произошло излишнее срабатывание БАВР.}{двигателя}}»}}
    \label{fig:ch1_img23}
\end{figure}

% DISS-005-B0565 .. DISS-005-B0566 (DISS-005-FIG-021 / DISS-005-TAB-015)
\begin{figure}[htbp]
    \centering
    \includegraphics[width=0.85\textwidth]{image24.png}
    \caption{Рисунок 2. Поведение РНМ.}
    \label{fig:ch1_img24}
\end{figure}

% DISS-005-B0567
\subsubsection{Рекомендации по наладке}

% DISS-005-B0568
При подключении мощной двигательной нагрузки рекомендуется уменьшать коэффициент отстройки с стандартного значения до 0,6. Более точное значение коэффициента можно определить по итогам пуска двигателя, в данном случае это можно осуществить по следующему алгоритму:

% DISS-005-B0569 .. DISS-005-B0571
\begin{itemize}[leftmargin=*,itemsep=0pt]
    \item вывести ключ БАВР
    \item Пустить двигатель
    \item По осциллограмме* оценить максимальный уровень колебания и его, с некоторой отстрой в сторону уменьшения, задать.
\end{itemize}

% DISS-005-B0572
*Примечание: Осциллограмма запишется за счёт встроенной логики функционирования пуска осциллографа при выведенном ключе БАВР. Так же можно временно ввести пуск осциллографа по защите 27/27S (ЗМН) или другим сигналам, которые считают целесообразными наладчики.

% DISS-005-B0574
\subsubsection{Распространённость проблемы}

% DISS-005-B0575
Проблемы необходимости ручной подстройки оптимального значения коэффициента отстройки характерно не для какого-то одного объекта, но для различной с мощной двигательной нагрузкой. Причём это значение зависит от множества факторов, влияющих на скорость уменьшения амплитуды тока. Где-то эта отстройка требуется большой, ввиду дополнительных колебаний амплитуды тока, как можем видеть на данном примере. А где-то, достаточно небольшой значения -- порядка 0,7-0,75, как можем видеть на другом примере (рис. 3).

% DISS-005-B0576 .. DISS-005-B0577 (DISS-005-FIG-022 / DISS-005-TAB-016)
\begin{figure}[htbp]
    \centering
    \includegraphics[width=0.85\textwidth]{image25.png}
    \caption[Рисунок 3. Характер изменения амплитуды тока на иной двигательной нагрузке]{Рисунок 3. Характер изменения амплитуды тока на иной двигательной \commentnote{DISS-005-C0060}{Евдаков Алексей}{Это\newline ЗРУ-6/6,3кВ Шахта Сибирская}{нагрузке}.}
    \label{fig:ch1_img25}
\end{figure}

% DISS-005-B0578
\subsubsection{Иные варианты срабатывания}

% DISS-005-B0579
Указанное отпадание РНМ может быть обусловлено не только фактом пуска мощного двигателя на указанной подстанции, но и иными переключениями во внешней сети, ведь с повышением чувствительность органов, повышается не только быстродействие, но и доля возможных излишних отпадание органа ввиду особенностей расчёта векторов на основе Фурье преобразований.

% DISS-005-B0580
Так, например, случилось на следующем приведённом \commentnote{DISS-005-C0061}{Евдаков Алексей}{на ГПП-3 на заводе ТАНЕКО}{примере. }Здесь при проведении крупных ремонтных работ во внешней сети и существенных переключениях вызвавшим изменение угла между секций с 2 до 13 градусов. Ввиду данного переходного процесса, произошло колебание всех векторов и срабатывание БАВР.

% DISS-005-B0581 .. DISS-005-B0582 (DISS-005-FIG-023 / DISS-005-TAB-017)
\begin{figure}[htbp]
    \centering
    \includegraphics[width=0.85\textwidth]{image26.png}
    \caption{Рисунок 4. Поведение РНМ при переключениях во внешней сети на ТАНЕКО (ГПП-3).}
    \label{fig:ch1_img26}
\end{figure}

% DISS-005-B0584
Отметим, что в срабатывание произошло по углу, и данное функционирование считается корректным, так как угол равный 13 градусов выходит за область нормальных. Однако в ряде сетей это является распространённым явлением при проведении летних ремонтных работ.

% DISS-005-B0585
В данном случае был доработан орган пуска по углу с целями исключения ложного срабатывания. Но отметим, что улучшение функционирования устройства БАВР в целом, за счёт других органов, не закрывает вопроса оптимизации реле направления мощности, в частности\commentnote{DISS-005-C0062}{Feel}{(}{. Так, например, поиск оптимального уровня коэффициента отстройки для конкретного объекта является сложной задачей, если и вообще решаемой. }
